\section{From potentials to maps}\label{sec:gradient}

We first recover gradient control from the potential estimate.  The argument
is a coordinate-slice version of the convex-function interpolation inequality
of Delalande and M\'erigot \cite[Proposition~4.1]{DelalandeMerigot2023}.
Their inequality already applies to every convex body considered here.
We give a coordinate proof with explicit constants, using only convexity
and the volumes of coordinate projections.

\begin{lemma}[Gradient interpolation]\label{gradient:interpolation}
Let $u_0,u_1$ be convex, $L$-Lipschitz functions on $K$, and put
$f=u_0-u_1$.  For the coordinate vectors $e_j$, set
\[
 A_j=\bigl|\operatorname{proj}_{e_j^\perp}K\bigr|_{d-1},
 \qquad P_K=\sum_{j=1}^d A_j.
\]
For every $h>0$,
\begin{equation}\label{gradient:slice-bound}
 \|\nabla u_0-\nabla u_1\|_{L^2(\rho)}^2
 \le \frac{12d}{h^2}\|f\|_{L^2(\rho)}^2
       +\frac{28L^2P_K}{|K|}h.
\end{equation}
Only the almost everywhere coordinate derivatives are needed in this
inequality.
\end{lemma}

\begin{proof}
Fix $j$ and write
\[
 K_{j,h}=\{x\in K:x+he_j\in K\},\qquad
 D_{j,h}u(x)=\frac{u(x+he_j)-u(x)}{h}
 \quad(x\in K_{j,h}).
\]
Almost every nonempty coordinate slice of $K$ is an interval
$[\alpha,\beta]$ on which the restriction of each $u_i$ is convex and
$L$-Lipschitz.  Its derivative exists almost everywhere, is monotone, and
has absolute value at most $L$.  Fubini therefore gives all coordinate
derivatives almost everywhere in $K$.  For either function $u$, convexity
gives, at points where both derivatives exist,
\[
 0\le \delta_u(x):=D_{j,h}u(x)-\partial_j u(x)
 \le \partial_j u(x+he_j)-\partial_j u(x)\le 2L.
\]
If $\beta-\alpha\ge h$, cancellation of translated integrals gives
\begin{align*}
 &\int_\alpha^{\beta-h}
       \bigl(\partial_j u(t+h)-\partial_j u(t)\bigr)\,dt\\
 &\hspace{15mm}=
   \int_{\beta-h}^{\beta}\partial_j u(t)\,dt
       -\int_\alpha^{\alpha+h}\partial_j u(t)\,dt
   \le 2Lh.
\end{align*}
The coordinates transverse to the slice have been suppressed here.
If the slice length is less than $h$, its allowed portion is empty.
Thus, in all cases, integrating $\delta_u^2\le 2L\delta_u$ over slices
and then over the orthogonal projection yields
\begin{equation}\label{gradient:quotient-error}
 \int_{K_{j,h}}\delta_u^2\,d\rho
 \le \frac{4L^2hA_j}{|K|}.
\end{equation}
The excluded portion of any slice has length at most $h$, so
\[
 \rho(K\setminus K_{j,h})\le \frac{hA_j}{|K|}.
\]
There we use $|\partial_j f|\le 2L$.

On the allowed set, both $x$ and $x+he_j$ belong to $K$.  Translation
invariance of Lebesgue measure and $(a-b)^2\le 2a^2+2b^2$ give
\[
 \int_{K_{j,h}}|D_{j,h}f|^2\,d\rho
 \le \frac{4}{h^2}\|f\|_{L^2(\rho)}^2.
\]
Since
$\partial_jf=D_{j,h}f-\delta_{u_0}+\delta_{u_1}$, the inequality
$(a+b+c)^2\le 3(a^2+b^2+c^2)$ and
\eqref{gradient:quotient-error} imply
\[
 \int_K|\partial_jf|^2\,d\rho
 \le \frac{12}{h^2}\|f\|_{L^2(\rho)}^2
       +\frac{(24+4)L^2hA_j}{|K|}.
\]
Summing over $j$ proves \eqref{gradient:slice-bound}.  All projection
measures are finite because $K$ is compact.  The slice argument applies
to every $h>0$, including values larger than some or all slice lengths.

The coordinate derivatives form the full gradient almost everywhere.
Indeed, at an interior point $x$ where all partials of a convex function
$u$ exist, set $p_j=\partial_j u(x)$ and $q=\|v\|_1$ for a small nonzero
vector $v$.  The point $x+v$ is the convex combination of
$x+\operatorname{sign}(v_j)q e_j$ with weights $|v_j|/q$, omitting zero
weights.  The finitely many coordinate derivative limits and convexity
give $u(x+v)-u(x)\le p\cdot v+o(q)$.  Apply this bound to $-v$ and use
midpoint convexity to obtain the reverse bound.  Since
$q\le\sqrt d\,|v|$, the function is differentiable at $x$ with gradient
$p$.  The boundary of $K$ has zero volume, as proved in
\Cref{cells:calculus}, so this applies almost everywhere to both $u_i$.
\end{proof}

\begin{corollary}[Finite target stability]\label{gradient:finite}
Set $C_0=162$ and
\begin{equation}\label{gradient:constant}
 C_*=
 \left(12dR^2(1+\sqrt{C_0})^2
             +\frac{28L^2P_K}{|K|}\right)^{1/2}.
\end{equation}
For any two finitely supported probabilities $\mu,\nu$ on $Y$,
\[
 \|T_\mu-T_\nu\|_{L^2(\rho)}
 \le C_*W_2(\mu,\nu)^{1/3}.
\]
\end{corollary}

\begin{proof}
Write $w=W_2(\mu,\nu)$.  By \Cref{finite:potentials}, the maps have
convex, $L$-Lipschitz maximum potentials $u_0,u_1$ satisfying
\[
 \|u_0-u_1\|_{L^2(\rho)}
 \le (1+\sqrt{C_0})Rw.
\]
If $w>0$, apply \Cref{gradient:interpolation} with $h=w^{2/3}$.
The resulting bound on the squared map norm is $C_*^2w^{2/3}$.
Taking its square root proves the claim.  If $w=0$, the potential
difference has zero $L^2(\rho)$ norm, and sending $h$ to zero in
\eqref{gradient:slice-bound} gives the same conclusion.
\end{proof}

\begin{proposition}[Approximation and optimal maps]\label{gradient:extension}
For every $\lambda\in\Prob(Y)$ there exists a unique quadratic optimal
map $T_\lambda$ from $\rho$ to $\lambda$, up to a $\rho$-null set.
If $F_n\subset Y$ are finite $\epsilon_n$-nets with
$\epsilon_n\to0$, and $Q_n:Y\to F_n$ are Borel nearest-point
selections, then
\[
 T_{(Q_n)_\#\lambda}\longrightarrow T_\lambda
 \quad\hbox{in }L^2(\rho).
\]
\end{proposition}

\begin{proof}
Compactness of $Y$ gives the finite nets.  Enumerating each net and
choosing the first closest point gives a Borel selection with
$|Q_n(y)-y|\le\epsilon_n$.  Put $\lambda_n=(Q_n)_\#\lambda$,
discard any zero-mass sites, and let $T_n=T_{\lambda_n}$.
These finite maps exist by \Cref{cells:weights,cells:optimality}.
The coupling induced by $y\mapsto(Q_n(y),Q_m(y))$ gives
\[
 W_2(\lambda_n,\lambda_m)\le\epsilon_n+\epsilon_m.
\]
By \Cref{gradient:finite}, $(T_n)$ is Cauchy in $L^2(\rho)$.
Let $S$ be its limit and choose a subsequence converging almost
everywhere.  Its limit lies in $Y$ almost everywhere because $Y$ is
closed.  Modifying $S$ on a null set gives a Borel map with values in
$Y$.  For any continuous $\phi:Y\to\R$, bounded convergence and
uniform continuity give
\[
 \int_K\phi(S(x))\,d\rho(x)
 =\lim_n\int_Y\phi(Q_n(y))\,d\lambda(y)
 =\int_Y\phi(y)\,d\lambda(y).
\]
Here and below a limit along the chosen subsequence is understood where
needed.  It follows that $S_\#\rho=\lambda$.

To prove optimality directly, choose finite maximum potentials $u_n$
for $T_n$ and subtract a constant so that $u_n(x_*)=0$ at a fixed
$x_*\in K$.  Their slopes lie in $Y$, so they are $L$-Lipschitz and
$|u_n|\le L\diam(K)$ on $K$.  Arzel\`a--Ascoli gives a further
subsequence converging uniformly on $K$ to a convex, $L$-Lipschitz
function $u$.  Intersect the full-measure sets on which each of the
countably many finite-map supporting inequalities holds with the set
on which the selected subsequence of $T_n$ converges.  On this common set,
\[
 u_n(z)\ge u_n(x)+T_n(x)\cdot(z-x)\qquad(z\in K).
\]
Passing to the limit gives, simultaneously for every $z\in K$,
\begin{equation}\label{gradient:limit-support}
 u(z)\ge u(x)+S(x)\cdot(z-x)
 \quad\hbox{for $\rho$-almost every }x.
\end{equation}
Define the restricted convex conjugate
\[
 u^*(y)=\max_{z\in K}\bigl(z\cdot y-u(z)\bigr).
\]
It is finite and continuous, and hence bounded on $Y$.  The defining
inequality and \eqref{gradient:limit-support} give
\begin{equation}\label{gradient:fenchel}
 x\cdot y\le u(x)+u^*(y)\quad(x\in K,\ y\in Y),
 \qquad
 x\cdot S(x)=u(x)+u^*(S(x))\quad\rho\hbox{-almost everywhere}.
\end{equation}
For any competing map $V$ with $V_\#\rho=\lambda$, the integral of
$u(x)+u^*(V(x))$ is fixed by the marginals.  Thus
\eqref{gradient:fenchel} shows that $S$ maximizes
$\int_Kx\cdot V(x)\,d\rho(x)$.  The identity
\[
 \int_K|x-V(x)|^2\,d\rho(x)
 =\int_K|x|^2\,d\rho(x)+\int_Y|y|^2\,d\lambda(y)
       -2\int_Kx\cdot V(x)\,d\rho(x)
\]
proves quadratic optimality.

If $V$ has the same optimal cost, its nonnegative Fenchel gap
$u(x)+u^*(V(x))-x\cdot V(x)$ has integral zero. Equality and the
definition of $u^*$ imply
$u(z)\ge u(x)+V(x)\cdot(z-x)$ for every $z\in K$ at almost every $x$.
On almost every interior point of $K$, all coordinate
partials of $u$ exist: on each coordinate interval the one-sided
derivatives of a convex function are monotone, and Fubini applies.
At such a point $x$, a supporting slope $p$ must satisfy
\[
 p_j\le\lim_{t\downarrow0}\frac{u(x+te_j)-u(x)}{t},
 \qquad
 p_j\ge\lim_{t\uparrow0}\frac{u(x+te_j)-u(x)}{t}.
\]
Both limits equal $\partial_j u(x)$, so the supporting slope is
unique.  The differentiability argument in the proof of
\Cref{gradient:interpolation} identifies it with $\nabla u(x)$.
The boundary of $K$ has zero volume by \Cref{cells:calculus}; hence
$V=S$ almost everywhere.
The same integrated inequality also proves optimality among arbitrary
transport plans.  Equality forces their target coordinate to be this
unique supporting slope over almost every source point, so the optimal
plan is the one induced by $S$.

We have identified the $L^2$ limit as the unique map $T_\lambda$.
The original sequence was Cauchy, so the asserted convergence holds
for the full sequence.
\end{proof}

\begin{proof}[Proof of \Cref{main:stability}]
Use the same finite nets and selections for $\mu$ and $\nu$, and set
$\mu_n=(Q_n)_\#\mu$, $\nu_n=(Q_n)_\#\nu$.  For any coupling
$\pi$ of $\mu,\nu$, its image under $(Q_n,Q_n)$ is a coupling of
$\mu_n,\nu_n$.  The triangle inequality in $L^2(\pi)$ gives
\[
 W_2(\mu_n,\nu_n)
 \le\left(\int_{Y\times Y}|y-z|^2\,d\pi(y,z)\right)^{1/2}
       +2\epsilon_n.
\]
Taking the infimum over $\pi$ yields
$W_2(\mu_n,\nu_n)\le W_2(\mu,\nu)+2\epsilon_n$.
Apply \Cref{gradient:finite} and then use the strong $L^2(\rho)$
convergence from \Cref{gradient:extension} to conclude
\[
 \|T_\mu-T_\nu\|_{L^2(\rho)}
 \le C_*W_2(\mu,\nu)^{1/3}.
\]
The constant in \eqref{gradient:constant} depends only on $K$ and an
enclosing radius for $Y$.  In particular it is independent of atom
counts, masses, site separations, or target densities.  No convexity
of $Y$ is required: the site interpolation used earlier stays in the
ball of radius $L$, and the discretizations take their values in $Y$.
The argument includes a singleton $Y$, for which all target maps
coincide.
The sharpness assertion was proved in \Cref{sharpness:exponent}.
\end{proof}
